\subsection{Introduction to Derivatives}

\begin{definition}[Derivative]
The derivative measures the instantaneous rate of change of a function with respect to its input. Geometrically, it is the slope of the tangent line to the graph at a point.

Common notations are
\[
  f'(x) = \lim_{h \to 0} \frac{f(x+h)-f(x)}{h},
\]
\[
  m_T = \lim_{h \to 0} \frac{f(x+h)-f(x)}{h},
\]
\[
  \frac{dy}{dx} = \lim_{h \to 0} \frac{f(x+h)-f(x)}{h}.
\]
\end{definition}

\noindent Useful notation:
\begin{itemize}
  \item Average rate of change over an interval: $\dfrac{\Delta y}{\Delta x}$.
  \item Instantaneous rate of change at a point: $\dfrac{dy}{dx}$.
\end{itemize}

For a function $f(x)$ at $x=a$,
\[
\left.\frac{dy}{dx}\right|_{x=a} = \lim_{h\to 0}\frac{f(a+h)-f(a)}{h}
\]
and equivalently
\[
f'(a) = \lim_{h\to 0}\frac{f(a+h)-f(a)}{h}.
\]

Another equivalent form is
\[
\lim_{x \to a}\frac{f(x)-f(a)}{x-a}.
\]

\begin{example}
Given $f(x)=\dfrac{1}{x}$, find the slope of the tangent at $x=2$.
\begin{align*}
m &= \lim_{x\to 2}\frac{\frac{1}{x}-\frac{1}{2}}{x-2} \\
  &= \lim_{x\to 2}\frac{2-x}{2x(x-2)} \\
  &= \lim_{x\to 2}\left(-\frac{1}{2x}\right) \\
  &= -\frac{1}{4}.
\end{align*}
\end{example}

The derivative may fail to exist at $x=a$ when:
\begin{itemize}
  \item the function is discontinuous at $x=a$;
  \item the function has a corner at $x=a$;
  \item the function has a cusp at $x=a$;
  \item the tangent is vertical at $x=a$.
\end{itemize}
